\section{Data Records}

\chinatherm{} contains 25 annual LST products for 2000--2024, each comprising 91 tiles of 4° by 4°. The record consists of 2,275 single-band GeoTIFF files with a total size of approximately 611.03~GiB. The GeoTIFFs use Deflate compression. ScienceDB provides 25 annual TAR archives, together with an English README and data dictionary\supercite{liang2026chinatherm30}.

The GeoTIFFs use the EPSG:4326 coordinate reference system and an angular pixel spacing of approximately 0.0002695°. Values are stored as Int16, with a scale factor of 0.01~°C and a NoData value of $-32768$. Filenames identify the product year and the integer longitude and latitude of the tile's western and southern edges. Field definitions are provided in the accompanying data dictionary.

The scientific domain is land within China. The core files retain rectangular tiles for mosaicking, so water areas and portions of boundary tiles outside China may contain model outputs; land statistics require the corresponding national and land boundaries. Figure~\ref*{fig:data-records} shows the national distributions for 2000 and 2024 and progressively finer views of western Sichuan and Beijing, from 4° tiles to local 30 m patterns.

\begin{figure}[H]
\centering
\figureorplaceholder{../figures/main/Figure3.pdf}{}{7.4cm}
\ReviewFloatBegin
\caption{National coverage and multiscale spatial views of \chinatherm{} in 2000 and 2024. a, h, National overviews at 0.01° for 2000 and 2024, respectively. b--d and e--g, The 4° tiles, 0.75° local windows and 0.26° close-up windows for western Sichuan (X100\_Y30) and Beijing--Yanshan (X116\_Y38) in 2000. i--k and l--n, Corresponding views at the same locations and scales in 2024.}
\label{fig:data-records}
\ReviewFloatEnd
\end{figure}
